\section*{Chapter 11 --- Testing without hardware}
\addcontentsline{toc}{section}{Chapter 11 --- Testing without hardware}

\solution{ex:11:1}{Get the suite green}
The exercise asks for code. The log, on the other hand, is the point: write down what the test
name said, where the bug was, and how long it took.

What you usually see is that the time falls sharply as the exercise goes on, and the reason is not
that the bugs get easier but that you stop guessing. The procedure in \secref{sec:debugging} ---
test name, index, expected value, and \emph{only then} your own code --- is the whole difference.

\solution{ex:11:2}{Which double catches what?}
\ans{a} \textbf{Scripted.} The recorded bytes end up in the wrong transaction, which shows
immediately. \code{BankTransport} obediently writes the values into the registers the command byte
points out and notices nothing.
\ans{b} \textbf{Both.} Scripted sees that the sent byte is \code{01} where the test expects
\code{00}; Bank sees that the error bit never goes low.
\ans{c} \textbf{Bank only.} From the scripted transport's point of view exactly the right bytes
were sent --- the fault is in what the driver hands back in \code{frame.data}, not in what it
sent.
\ans{d} \textbf{Both.} Scripted is missing the \code{RX_ACK} transaction in the recording; Bank
sees that RX valid stays set and that the next \code{receive()} gives the same frame.
\ans{e} \textbf{Bank only}, indirectly. Scripted checks what was \emph{sent}, and that byte is
sent correctly --- the fault is that the answer is assembled wrongly. Bank gives back values that
do not match what was written, so some read comes out wrong.
\ans{f} \textbf{Both.} Scripted sees that transactions go out even though the frame is invalid;
Bank sees that the registers are written.
\ans{g} \textbf{Neither.} And that is precisely why the row is there: a wiring fault is invisible
to every test run on the laptop, however good the test double is. That is what \kapref{ch:bringup}
exists for, and the reason rung 1 of the bring-up ladder is a loopback.

\solution{ex:11:3}{Read the suite as a requirements specification}
\ans{a} The test case checking the \code{TX_DATA_LO} transaction on a \code{send()} of
\code{id = 0x123}, \code{dlc = 2}, data \code{AA BB}. Its five \code{EXPECT_EQ} lines are the
example's five bytes, written as a requirement.
\ans{b} The case sending a frame with \code{dlc = 9} and checking that the recorded byte stream is
\emph{empty}, that is, that no transactions went out at all. It checks not only the return value
but that the validation happens \emph{first}.
\ans{c} The distribution should lean towards \code{BankTransport} for behaviour and
\code{ScriptedTransport} for the protocol --- roughly as in the table in \secref{sec:twodoubles}.
If you see only one of them in the suite, a whole class of bugs is unaccounted for.

\solution{ex:11:4}{The whole stack}
The exercise asks you to run a test case; the questions are answered here.
\ans{a} Zero. \code{EchoNode} is unchanged since \kapref{ch:komponent}.
\ans{b} That \code{EchoNode} depends on \code{driver::can::Interface} and nothing else. Both
\code{Stub} and \code{Spi} fulfil that contract, so the swap is a single line in the test's
Arrange part.
\ans{c} Everything. \code{EchoNode} would have had to include \code{spi.h}, and thereby dragged in
\code{ByteTransport} and the whole transport layer. The test would have needed a transport just to
construct the component, and the component test in \kapref{ch:komponent} --- which runs entirely
without a transport --- could not have been written.

\solution{ex:11:5}{What remains}
At least these, and the last one is the one usually forgotten:

\begin{itemize}
  \item That the \textbf{interpretation} of the contract is right. The suite is written by a human
    who read the same document as you.
  \item That the \textbf{hardware group} interpreted it the same way. Their testbenches are
    written against the same document, by somebody else.
  \item That the \textbf{wiring} is right --- a swapped \code{MOSI}/\code{MISO} is invisible here.
  \item That the \textbf{time budget} holds. \code{BankTransport} answers immediately.
  \item That you \textbf{poll often enough} not to lose frames.
  \item That \code{AvrSpiTransport} works. It is the only class in the project that no automated
    test covers, since it talks to a physical register.
\end{itemize}
